\section{Binary boundary tensors for faithful Jones computation}
\label{sec:spin-jones}

\subsection{The two incoming reports and the integration decision}
Commit \code{1be2abc8c} supplies two research archives dated 8 October 2026.
The binary-tensor report develops a complete Jones evaluator, a near-prefix
compressed matcher, and dihedral boundary transport. The modular-boundary
report develops finite-field reuse of the marked Euler observer. Both target
upstream \code{8a9583494}; neither snapshot includes our subsequent window,
phase, sparse-endpoint, or endpoint-progression refinements. Their patches
must therefore be integrated by component, preserving those later changes.
The supplied manifests verify all 480 binary-package files and all 380
modular-package files. These integrity checks establish source identity,
not mathematical correctness.

This section integrates the binary Jones component, its three test files,
benchmark driver and historical policy sources. The other components remain
under review. In particular, the modular report reports improved repeated
boundary queries but slower whole raw scans on all seven measured examples;
its fixed prime palette can be inconclusive. The new near-prefix branch
must be compared with the current matcher, since its archived baseline
predates our window optimization. No speedup from either unintegrated
component is attributed to the maintained recognizer here.

The new optional \code{spin-faithful} backend improves the maintained full
Jones bound from the faithful Potts implementation's
$\operatorname{poly}(n)2^{O(\sqrt n\log n)}$ to
$\operatorname{poly}(n)2^{O(\sqrt n)}$. Comparable general bounds are known
in the literature, including Maria~\cite{quantum-width}; this is an
implementation improvement, not a claim of a new complexity class result.
Jones polynomial one remains inconclusive for unknot recognition.

\subsection{An integral turning certificate from the PD rotation system}
For a validated nonempty classical knot diagram, let $\alpha$ exchange the
two darts on an edge and let $\sigma$ advance one slot counterclockwise at
a crossing. The cycles of $\phi=\sigma\alpha$ are the $n+2$ faces. Pick an
outer face $f_\infty$. An antisymmetric integral cochain $b$ on directed edges
must satisfy
\[
 b_{\alpha(d)}=-b_d,\qquad
 \sum_{d\in f}b_d=|f|-4+8\mathbf1_{f=f_\infty}.
\]
A face corner turns by $-1$ quarter-turn, so these are the edge contributions
to the total turning $-4$ of a bounded face and $+4$ of the outer face.
The demands sum to zero, since the face lengths sum to $4n$.

Root a dual spanning tree at the outer face. On each child face's parent
edge, assign the total demand of its subtree, with opposite sign on the
reverse dart; assign zero to the other edges. At each face, the parent
subtree total minus its children's totals equals its demand. The root
follows from total demand zero. Each edge value has magnitude at most
$4n+4$, and the sum $L_b=\sum_e|b_e|$ is at most $4(n+1)^2$.
The certificate therefore has $O(n)$ integers of $O(\log n)$ bits.
An independent verifier reconstructs the PD faces and checks the displayed
equations directly, without repeating the tree solver.

Why do these equations suffice for all smoothing circles? In any smooth
planar realization with cardinal crossing tangents, geometric edge turns
satisfy the same equations. The difference from $b$ integrates to zero
around every face. On the sphere, face boundaries span the cycle space,
so that difference integrates to zero around every closed graph walk.
Replacing geometric edge turns by $b$ therefore leaves each oriented
smoothing circle's total turning unchanged: it is $+4$ or $-4$.
This proof permits loops, repeated face vertices, and a smoothing circle
visiting the same crossing twice. It does not apply unchanged to virtual
diagrams on surfaces of positive genus.

\subsection{Two orientations replace boundary connectivity}
Write the unnormalized bracket as
\[
 \mathcal B_D(A)=\sum_{s\in\{0,1\}^n}A^{n-2|s|}\delta^{c(s)},
 \qquad \delta=-A^2-A^{-2}.
\]
Introduce $\rho$ with $\rho^4=-A^2$. Orient every smoothing circle.
Its two orientations contribute $\rho^4+\rho^{-4}=\delta$ by the turning
certificate. Thus the same state sum can be expressed with one binary
arrow variable on each diagram edge. At a crossing, retain only smoothings
whose pairs have one incoming and one outgoing arrow; multiply by their
usual $A$ weight and by $\rho$ to the sum of their local turns. On an edge,
multiply by $\rho^{\epsilon b_e}$, where $\epsilon$ is its chosen arrow.
Expanding this tensor product gives precisely all oriented smoothing
states and recovers the bracket.

A crossing order with $w$ cut edges now has at most $2^w$ tensor keys.
There is no connected-prefix or common-disk assumption: a key records
arrows, not a pairing of boundary points. The local tensor has bounded
arity, and self-loop arrows are summed at their crossing. This exact
summary preserves the bracket under every remaining tensor contraction.
It does not preserve the relative Khovanov complex; cancellations in a
polynomial may erase data needed for homology.

\subsection{One phase and one integer per key}
Set $A=B^2$ and $\rho=\zeta B$, where $\zeta^4=-1$. Initially this suggests
four integer coordinates per tensor key. In fact the phase modulo four
is determined by the boundary arrows. To see this, for an edge from slot
$j$ at $v$ to slot $k$ at $u$, the cochain
$b_d+j-k-2$ has zero face integrals modulo four. It is therefore a vertex
potential difference $\nu(u)-\nu(v)$. At a crossing the local turning is
$\sum_d\epsilon_d j$ modulo four, and $\sum_d\epsilon_d=0$.
After adding the vanishing potential terms, each internal edge contributes
$2\epsilon_d=2$ modulo four, independent of its arrow. The remaining terms
are determined by the boundary. Hence one phase $r\in\{0,1,2,3\}$ and one
signed integer suffice. The implementation checks phase agreement whenever
terms merge, and checks phase zero for the final scalar.

Multiply each crossing tensor by $B^4$ and each edge tensor by $B^{|b_e|}$.
Every local exponent is then nonnegative. Local multiplication requires
only shifts and additions of signed integers, with phase/sign updates.
If $K_0=4n+L_b$ and $K=\max(K_0,6n+4)$, the final padded scalar is
\[
 S=B^K\mathcal B_D(B^2).
\]
All expanded shifted monomials have $B$-degree at most $8n+2L_b$ before
final padding. There are at most $2^{3n}$ smoothing--arrow assignments.
The triangle inequality consequently bounds every intermediate integer
by $O(n^3)$ bits for the base chosen next; large cancellation does not
invalidate this upper bound.

\subsection{Faithful evaluation and coefficient recovery}
For a knot with $n$ crossings, the normalized Jones polynomial
$V_D(t)=\sum_jc_jt^j$ obeys
\[
 \sum_j|c_j|\le4^n,\qquad |j|\le2n\quad(c_j\ne0).
\]
The first follows from $2^n$ states, each with at most $n+1$ circles and
coefficient norm at most $2^n$ after reduction. Smoothing degrees and the
writhe correction give the second, deliberately loose, degree bound.
Choose
\[
 \beta=\left\lceil(2n+2)/8\right\rceil,\quad B=2^\beta,\quad X=B^8.
\]
Then $X\ge4\cdot4^n$. Balanced base-$X$ digits uniquely recover every
coefficient after multiplying the Laurent polynomial by a known monomial.
This is an injective exact encoding, not a random evaluation.

If $a$ is the writhe, the expected scaled unknot bracket is
\[
 U=(-1)^{a+1}(B^8+1)B^{K+6a-4}.
\]
Its exponent is nonnegative. Equality $S=U$ proves $V_D=1$; inequality
proves knottedness. For the full polynomial, form the exact integer
\[
 X^{2n}V_D(X^{-1})
   =\frac{(-1)^{a+1}S B^{16n+4-K-6a}}{X+1},
\]
moving a negative power to the denominator. Check divisibility, then
recover balanced digits $d_i$ and set $c_j=d_{2n-j}$. The decoder checks
normalization, degree and coefficient bounds, including before its identity
shortcut. These checks validate the decoding interface; they are not an
independent certificate of a supplied scalar's contraction. The empty PD
represents one unknot circle and receives its explicit special case.

\subsection{Complexity and recognition semantics}
Each crossing has a bounded number of transitions per binary frontier key,
and all keys and integer coefficients have polynomial bit length. A fixed
width-$w$ order therefore costs $\operatorname{poly}(n)2^{O(w)}$
deterministic bit operations. Even charging linear scans for worst-case
Python dictionary collisions only squares the exponential factor and
preserves this bound. Our existing verified separator order has
$w=O(\sqrt n)$. The default spin query considers a greedy candidate even
when the caller supplies an order, and retains an order within the
separator certificate's bound. With local caps and deadlines disabled,
full Jones computation thus has the stated square-root exponential bound.
For supplied orders of width $O(\log^2 n)$ the corresponding fixed-order
query is quasi-polynomial.

The production default remains unchanged: the report exhibits regressions
as well as improvements. Users can select \code{spin-faithful} in either
the Jones command or the recognizer's Jones-filter option. Local state or
transition exhaustion publishes no unfinished scalar or identity claim;
a completed separator certificate may still be reused. Caller cancellation
propagates. Large evidence integers are serialized in hexadecimal without
changing global decimal conversion settings. A polynomial-one result
continues to the independent exact recognizer. Consequently this backend
does not establish a general subexponential unknot-recognition algorithm.

\subsection{Independent validation and controlled experiments}
The imported tests perform 140 full-polynomial comparisons with the
independent whole-cube Laurent oracle, covering mirrors, crossing orders
and outer faces. They verify oriented-circle turning directly, test
20 vertex gauges, exercise 180 bounded signed polynomial decodes, reject
forged normalization, check budget/cancellation behavior, and verify CLI
and recognition fallback behavior. A shuffled 64-crossing grid guards
against the report's earlier supplied-order policy defect.

All 714 maintained tests pass in 164.627 seconds with Regina 7.4.1.
The integrated benchmark uses the delivered final source, including the
normalization check added after the report's historical timing runs.
It records 490 measurements and 70 excluded warmups over fourteen inputs,
five arms and seven shuffled rounds (seed 26100860). There are 469 complete
measurements and 21 local-cap outcomes. Default queries include fresh PD
validation and order preparation; the fixed-order scope uses the same
externally prepared certified order and includes all tensor and cochain setup.
Each query requests the full polynomial, with 4096 states, 200000 transitions
and a three-second cooperative deadline. Completed arms agree; small inputs
use an independent cube oracle, and grids have checked descending certificates.

\begin{center}\small
\begin{tabular}{@{}llrrrrr@{}}
\toprule
Input & $n$ & Potts & A/A & Spin & Fixed P & Fixed S \\
\midrule
Trefoil & 3 & 0.190 & 0.187 & 0.544 & 0.158 & 0.416 \\
Conway & 11 & 0.638 & 0.655 & 2.314 & 0.445 & 1.568 \\
KT & 11 & 0.667 & 0.647 & 2.348 & 0.491 & 1.576 \\
Hard 8 & 8 & 0.381 & 0.372 & 1.287 & 0.283 & 0.993 \\
Braid5 36 & 36 & 3.104 & 3.112 & 11.495 & 2.755 & 8.455 \\
Weaving 10 & 10 & 0.698 & 0.647 & 2.515 & 0.522 & 1.784 \\
Torus 101 & 101 & 37.243 & 37.376 & 39.350 & 34.529 & 32.174 \\
Tree 6 & 126 & 10.277 & 10.341 & 62.948 & 7.905 & 48.508 \\
Tree 7 & 254 & 73.257 & 73.270 & 394.441 & 68.170 & 352.804 \\
Grid 8 & 64 & 6.791 & 6.747 & 23.368 & 6.119 & 17.737 \\
Grid 10 & 100 & 39.177 & 39.008 & 67.193 & 39.033 & 57.578 \\
Grid 12 & 144 & 380.542 & 459.847 & 306.520 & 369.836 & 224.954 \\
Grid 14 & 196 & limit & limit & 732.041 & limit & 717.790 \\
Grid 8 shuffled & 64 & 29.152 & 29.297 & 22.882 & 6.043 & 17.528 \\
\bottomrule
\end{tabular}
\end{center}
Median milliseconds for full Jones queries. Fixed P/S share an externally prepared order; other columns include ordering. Limits are censored, not completed-query times.


On the 144-crossing grid, ordinary Potts/A/A/spin medians are
380.542/459.847/306.520 milliseconds. On the common order they are
369.836/224.954 milliseconds for Potts/spin. Median within-round ratios are
1.638 and 1.648 respectively; they need not equal ratios of the displayed
medians. Potts stores 877 peak states and performs 48386 transitions;
spin stores 240 and performs 32098. On the 196-crossing grid, every Potts
arm reaches the state cap, while all spin queries complete with 480 peak
states and 89122 transitions. This is a capacity result, with no
completed-time ratio against the capped arms.

Regressions remain substantial. The 254-crossing tree medial has ordinary
Potts/spin medians 73.257/394.441 milliseconds and fixed-order medians
68.170/352.804. The Potts cancellation pattern leaves only one state on
that input, while spin reaches 272. On the shuffled 64-crossing grid,
ordinary medians favor spin, 29.152/22.882 milliseconds, but fixed-order
medians favor Potts, 6.043/17.528. The ordinary gain on this case comes
from ordering rather than a faster contraction. These full invariant
queries do not measure completed recognition on a hard-knot corpus.

Raw maintained measurements are in
\path{results/spin_jones_integrated_20261008.json}, with per-file source hashes.
Its inherited \code{baseline\_commit} field identifies the report's integration
base \code{8a9583494}; the actual compared source is identified by the hashes.
The original \path{spin_jones_20261008.json}, initial-policy results and ordering
follow-up remain separate historical records. The report's previous performance
numbers are not relabeled as results of this integrated run.
